% chapters/ch8_exp_interference.tex --- 4.7 高干渉環境での性能評価実験
\section{高干渉環境での性能評価実験}
\label{sec:exp_interference}

\subsection{実験の目的}

AP近接による高干渉環境で，TailGuard（連続受信失敗時のフォールバック機構）の有効性を検証する．

\subsection{実験結果}

\tabref{tab:e2_ccs_tail_guard}に結果を示す．CCS単独では$P_{\mathrm{out}}(2\mathrm{s})=10.77\%$と悪化したが，TailGuard追加により1.28\%に改善し，QoS制約を満たした．

\begin{table}[tb]
  \centering
  \caption{高干渉環境での結果}
  \begin{tabular}{lrrrr}
    \toprule
    条件 & 平均電流[mA] & $TL_{p95}$[ms] & $P_{\mathrm{out}}(2\mathrm{s})$ \\
    \midrule
    100ms & 96.93 & 251.7 & 0.00\% \\
    CCS & 90.12 & 4116.6 & 10.77\% \\
    CCS+TailGuard & 92.01 & 944.5 & 1.28\% \\
    \bottomrule
  \end{tabular}
\end{table}

\clearpage
